\sectionTitle{Berufserfahrung}{\faSuitcase}
\begin{experiences}
  \experience
    {März 2026} {Research Engineer}{NEC Laboratories Europe (NLE)}{Heidelberg, Vollzeit}
    {Juni 2025} {
      \begin{itemize}
        \item{Forschung zu Langzeitgedächtnis in agentischen LLMs und kausalen Erklärungen durch Prototypen und Mechanismen.}
        \item{Entwicklung skalierbarer Echtzeit-LLM-Systeme für die Bearbeitung von Notrufen mit FastAPI und CI/CD.}
      \end{itemize}
    }
    {Langzeitgedächtnis, Kausale Prototypen, Systemdesign und Software-Releases, Betreuung von Abschlussarbeiten und Führung}
  \emptySeparator
  \experience
    {Mai 2025} {Research Associate}{NEC Laboratories Europe (NLE)}{Heidelberg, Vollzeit}
    {Januar 2024} {
      \begin{itemize}
        \item{Entwicklung interpretierbarer Deep-Learning-Varianten des Classification-by-Components-Ansatzes auf dem Stand der Technik.}
        \item{Entwicklung einer interpretierbaren Webanwendung für Gesundheitsassistenz mit GLVQ- und GTLVQ-Ansätzen.}
      \end{itemize}
    }
    {Prototypenbasiertes Lernen, Prototypenentwicklung, Softwareentwicklung}
  \emptySeparator
  \experience
    {September 2023} {Werkstudent}{Institute for Enterprise Systems}{Mannheim, Teilzeit}
    {April 2023} {
      \begin{itemize}
        \item{Entwicklung von Datenvorverarbeitung und Codeoptimierungen für die BERT-Systempipeline.}
        \item{Entwicklung neuartiger Architekturverbesserungen für Smart-Meter-Zeitreihendaten.}
      \end{itemize}
    }
    {Zeitreihendaten, Transformer, Datenverarbeitung}
  \emptySeparator
  \experience
    {Dezember 2023} {Werkstudent}{trinamiX GmbH}{Ludwigshafen, Teilzeit}
    {April 2023} {
      \begin{itemize}
        \item{Entwicklung von Datenverarbeitung und Codeoptimierungen für eine CNN-basierte Deep-Learning-Pipeline.}
        \item{Benchmarking der Transfer-Learning-Performance von EfficientNet mit biometrischen Daten dynamischer Eingabegröße.}
        % \item{Quantitativer Vergleich der Tiefenschätzung von Mediapipe- und InsightFace-Lösungen zur Landmarkenerkennung.}
      \end{itemize}
    }
    {PyTorch, Hydra, MLflow}
  \emptySeparator
  \experience
    {Februar 2023} {Werkstudent}{FARO Technologies}{Korntal-Münchingen, Teilzeit}
    {September 2022} {
      \begin{itemize}
        \item{Durchführung großskaliger, hochauflösender Punktwolkensegmentierung \textit{(10\textsuperscript{7} Punkte)} mit RandLA-Net.}
        \item{Feinabstimmung von RandLA-Net auf synthetisch erzeugten Daten und Benchmarking der Verbesserungen.}
        % \item{Benchmarking der Modellinferenz von RandLA-Net mit synthetisch erzeugten Daten.}
        % \item{Anwendung von Segmentierungsmodellen auf realistische, sehr hochauflösende Punktwolkenscans.}
      \end{itemize}
    }
    {PyTorch, Open3D-ML, Erzeugung synthetischer Daten}
  \emptySeparator
  \experience
    {August 2022} {Data-Science-Praktikant}{ZEW}{Mannheim, Vollzeit}
    {Juni 2022} {
      \begin{itemize}
        \item{Entwicklung von Web-Scrapern für mehrsprachige Daten und von Übersetzungs-Wrappern zur Synchronisierung von Datensätzen.}
        \item{Statistische Analyse von Fußballkompetenzen mit Difference-in-Differences-Modellen.}
      \end{itemize}
    }
    {Selenium, Beautiful Soup, HuggingFace, PyTorch}
  \emptySeparator
  \experience
    {September 2021} {Data Engineering Specialist}{GE Healthcare}{Bengaluru, Vollzeit}
    {August 2019} {
      \begin{itemize}
        \item{Entwicklung eines POC zur prädiktiven Ausreißererkennung für MRT-Modalitäten aus mehreren Sensordatenquellen.}
        \item{Entwicklung von Kennzahlen zur Maschinenauslastung, um die zusätzliche Belastung von Gesundheitssystemen durch COVID-19-Patienten zu messen.}
        \item{Entwicklung und Bereitstellung von Ingestionsprozeduren im Terabyte-Maßstab mit \textit{PL/pgSQL} auf AWS Redshift.}
      \end{itemize}
    }
    {Python, R, PL/pgSQL, AWS, Redshift}
  \emptySeparator
  \iffalse
  \experience
    {August 2014} {Studienlücke | Vorbereitung auf mehrere Auswahlprüfungen}{XXX-Test bestanden}{Bengaluru}
    {August 2013} {
      \begin{itemize}
        \item{Prüfung X mit Ergebnis Y abgelegt. Unter den besten X\% der Teilnehmenden.}
        \item{Prüfung X mit Ergebnis Y abgelegt. Unter den besten X\% der Teilnehmenden.}
      \end{itemize}
    }
    {Höhere Analysis, Differentialgleichungen, Lineare Algebra, Allgemeinwissen}
  \emptySeparator
  \fi
  \experience
    {Juni 2019} {Data-Engineering-Praktikant}{American Express}{Bengaluru, Vollzeit}
    {Januar 2019} {
      \begin{itemize}
        % \item{Entwicklung von Prometheus- und Grafana-basierten Monitoring- und Alerting-Werkzeugen für über 1000 Couchbase-Datenbank-VMs.}
        \item{Entwicklung Prometheus- und Grafana-basierter Monitoring- und Alerting-Werkzeuge für Couchbase-DB-VMs.}
        \item{Konzeption, Entwicklung und Bereitstellung paralleler VM-Patching-Skripte zur vierfachen Reduzierung der Ausfallzeit.}
      \end{itemize}
    }
    {Couchbase, Python, Grafana, Prometheus}
  \emptySeparator
  \experience
    {Juli 2018} {Software-Engineering-Praktikant}{GE Digital}{Bengaluru, Vollzeit}
    {Juni 2018} {
      \begin{itemize}
        \item{Ergänzung eines SAML-Re-Authentifizierungsmoduls und Migration der Enovia-PLM-Codebasis auf eine REST-Architektur.}
        % \item{Implementierung einer REST-basierten Architektur aus bestehendem SOAP-Legacy-Code.}
        \item{Durchführung einer Fallstudie zu Open-Source-Werkzeugen für statische Analyse auf der Enovia-Java-Codebasis.}
      \end{itemize}
    }
    {REST, Java, SAML}
  \emptySeparator
\end{experiences}
